\renewcommand{\Anno}{2026}

\renewcommand{\NumeroProblema}{3}

\renewcommand{\Testo}{
Un'urna trasparente contiene 101 palline, numerate con 101 interi positivi distinti, di cui 51 dispari. Due giocatori bendati, A e B, estraggono alternativamente una pallina dall'urna, iniziando da~A, fino ad esaurimento delle palline. Le palline estratte non vengono reintrodotte nell'urna. Al termine, il giocatore~A avrà 51 palline, e il giocatore~B ne avrà 50. Ogni giocatore somma i numeri presenti sulle sue palline, e vince il giocatore che ha la somma maggiore.

Il giocatore~A, onesto, resta bendato e sceglie sempre a caso la pallina da estrarre, mentre il giocatore~B, ritenendosi svantaggiato dal giocare per secondo, si toglie la benda e ogni volta sceglie la pallina da ``estrarre'' dopo aver guardato i numeri di quelle rimaste.

\begin{enumerate}
\renewcommand{\labelenumi}{(\alph{enumi})}

\item  Dimostrare che il gioco non può finire in parità.

\item Dimostrare che esiste una scelta dei numeri presenti sulle palline per cui il giocatore~A ha la certezza di vincere.

\item Determinare se è possibile scegliere i numeri presenti sulle palline in modo che la probabilità di vittoria del giocatore~B risulti maggiore del 99.99\%.

\item Determinare se è possibile scegliere i numeri presenti sulle palline in modo che entrambi i giocatori abbiano il 50\% di probabilità di vincere.

\end{enumerate}
}

\renewcommand{\Risposta}{
% SSSUP-BOZZA-AI: Risposta; autore=ChatGPT; data=2026-11-08; stato=DA-VERIFICARE
% SSSUP-PROVENIENZA: AI: rielaborazione e sviluppo di appunti forniti
\begin{enumerate}
\item[(a)] La somma totale è dispari, quindi il pareggio è impossibile.
\item[(b)] Basta scegliere i numeri $2503,2504,\ldots,2603$.
\item[(c)] Sì: con $1,\ldots,96,10001,\ldots,10005$ e B che prende sempre una pallina grande disponibile, la probabilità di vittoria di B è
\[
1-\frac{15}{101\cdot99\cdot97}=1-\frac{15}{969903}>99{,}99\%.
\]
\item[(d)] No per strategie deterministiche di B, e quindi no se B massimizza la propria probabilità di vittoria. Se invece si ammettono strategie randomizzate deliberatamente non ottimali, il $50\%$ è possibile: questa distinzione va precisata nell'interpretazione del testo.
\end{enumerate}
}

\renewcommand{\Aiutino}{
% SSSUP-BOZZA-AI: Aiutino; autore=ChatGPT; data=2026-11-08; stato=DA-VERIFICARE
% SSSUP-PROVENIENZA: AI: rielaborazione e sviluppo di appunti forniti
Usa la parità per (a), numeri tutti vicini a una grande costante per (b), e cinque numeri dominanti rispetto agli altri per (c). Per (d), fissata una strategia deterministica di B, conta le foglie equiprobabili dell'albero delle scelte di A: il loro numero è un prodotto di numeri dispari.
}

\renewcommand{\SoluzioneUfficiale}{}

\renewcommand{\SoluzioniAlternative}{}

\renewcommand{\Osservazioni}{}

\renewcommand{\Commenti}{}

% Aggiunta redazionale: la soluzione preesistente non viene modificata.
\providecommand{\SoluzioneRedazionale}{}
\renewcommand{\SoluzioneRedazionale}{
% SSSUP-BOZZA-AI: SoluzioneRedazionale; autore=ChatGPT; data=2026-11-08; stato=DA-VERIFICARE
% SSSUP-PROVENIENZA: AI: rielaborazione e sviluppo di appunti forniti
Rielaboriamo la bozza fornita, esplicitando al punto (d) il ruolo della strategia di B.

\paragraph{Punto (a).}
La somma di $51$ interi dispari e $50$ interi pari è dispari. Se le due somme finali fossero uguali, la somma totale sarebbe pari. Quindi non ci sono pareggi.

\paragraph{Punto (b).}
Prendiamo $N+1,\ldots,N+101$, con $N$ pari. Ci sono esattamente $51$ dispari. La somma dei $51$ numeri più piccoli meno la somma dei $50$ più grandi vale
\[
51N+\frac{51\cdot52}{2}-\left(50N+\frac{101\cdot102}{2}-\frac{51\cdot52}{2}\right)=N-2499.
\]
Per esempio $N=2502$ rende questa differenza positiva. A vince pertanto anche nella peggiore assegnazione possibile, qualunque sia la strategia di B.

\paragraph{Punto (c).}
Scegliamo $96$ numeri piccoli $1,\ldots,96$ e cinque numeri grandi $M+1,\ldots,M+5$, con $M=10000$. La parità richiesta è rispettata: $48+3=51$ dispari. Poiché $M>4656+15$, possedere la maggioranza dei cinque numeri grandi basta a vincere, indipendentemente dai numeri piccoli.

B prende un numero grande ogni volta che ne resta almeno uno. Perché A ne prenda tre, deve pescarne uno in ognuna delle sue prime tre mosse: se in una di queste prende un numero piccolo, B esaurisce abbastanza presto i grandi da impedirgli la maggioranza. La probabilità di vittoria di A è quindi esattamente
\[
\frac5{101}\cdot\frac3{99}\cdot\frac1{97}
=\frac{15}{969903}<\frac1{10000}.
\]
B vince con probabilità maggiore del $99{,}99\%$.

\paragraph{Punto (d).}
Fissiamo una strategia deterministica di B, che può dipendere da tutta la storia precedente. Le scelte casuali di A formano un albero con $101,99,97,\ldots,1$ possibilità a livelli successivi. Le
\[
D=101\cdot99\cdot97\cdots3\cdot1
\]
foglie hanno tutte probabilità $1/D$. La probabilità di vittoria di ciascun giocatore è quindi un intero diviso per $D$, e non può essere $1/2$, dato che $D$ è dispari.

Anche se si ammettono strategie randomizzate, una strategia ottimale per B può essere scelta deterministica: per induzione a ritroso nell'albero finito, a ogni sua decisione basta scegliere una mossa che massimizzi la probabilità di vittoria futura. Pertanto il valore ottimale non può essere $1/2$.

Se invece B è autorizzato a randomizzare senza voler massimizzare la vittoria, la conclusione cambia. Una strategia che gli fa perdere certamente è prendere sempre il numero minimo rimasto: ogni sua pallina è minore di quella che A prenderà subito dopo, e A possiede inoltre la prima pallina. Per i numeri del punto (c), mescolando questa strategia perdente con quella che vince con probabilità $p_B>1/2$, e scegliendo quest'ultima con probabilità $1/(2p_B)$, B ottiene probabilità esattamente $1/2$. Dunque la risposta ``no'' richiede la convenzione deterministica oppure l'interpretazione di gioco ottimale.
}
